% META: {"id": "TSPE-PROBA-CR-013", "chapitre": "TSPE-PROBABILITES", "type_objet": "cours", "capacites_codes": ["C6", "C7"], "sources_inspiration": [], "mode_creation": "ex_nihilo", "status": "approved"}

\section{Somme de variables aléatoires, linéarité de l'espérance}

% ============================================================
% STRATE 1 — Decomposition et linearite
% ============================================================

\propriete[Décomposition]{
Une variable aleatoire complexe peut souvent s'ecrire comme somme de variables aleatoires plus simples. Par exemple, une variable $X$ suivant $\mathcal{B}(n,p)$ peut s'ecrire $X = X_1+X_2+\cdots+X_n$, ou chaque $X_i$ vaut $1$ si la $i$-eme epreuve est un succes, $0$ sinon (variable de Bernoulli).
}

\propriete[Linéarité de l'espérance, admise]{
Pour toutes variables aleatoires $X,Y$ (finies) et tous reels $a,b$ :
\[
  E(aX+bY) = a\,E(X) + b\,E(Y).
\]
En particulier, $E(X_1+\cdots+X_n) = E(X_1)+\cdots+E(X_n)$, \textbf{meme si les $X_i$ ne sont pas independantes}.
}

\margeAppui{
Contrairement a la variance (voir plus loin), la linearite de l'esperance ne necessite \emph{aucune} hypothese d'independance.
}

\exemple{
On lance $3$ des equilibres. $X$ = somme des points obtenus. Calculer $E(X)$.

$X=X_1+X_2+X_3$ ou $X_i$ est le resultat du $i$-eme de. $E(X_i) = \dfrac{1+2+3+4+5+6}{6} = \dfrac{7}{2}$ pour chaque $i$.

Par linearite : $E(X) = E(X_1)+E(X_2)+E(X_3) = 3\times\dfrac{7}{2} = \dfrac{21}{2} = 10{,}5$.

% BEGIN-VERIFY
% from sympy import Rational
% Ei = Rational(1+2+3+4+5+6, 6)
% assert Ei == Rational(7,2)
% E = 3*Ei
% assert E == Rational(21,2)
% END-VERIFY
}

% ============================================================
% STRATE 2 — Erreurs frequentes
% ============================================================

\erreurFrequente{
\textbf{Croire que la linearite de l'esperance necessite l'independance.} C'est faux : $E(X+Y)=E(X)+E(Y)$ est toujours vraie, meme si $X$ et $Y$ sont dependantes. Cette hypothese sera necessaire pour la \emph{variance} uniquement.
}
